\section{A restricted model: ordered decision diagrams}
Ordered binary decision diagrams are classical representations
\cite{Bryant1986}. We use deterministic \emph{layered} diagrams:
each layer reads the next bit in a fixed order, every path reads
all bits, and the final state supplies the answer. Equivalent
residual computations may merge. This convention avoids ambiguity
about skipped levels and early terminals.

For \(a,b\in\bits^m\), define
\[
 g_m(a,b)=\bigwedge_{i=1}^m\neg(a_i\land b_i).
\]
This is the satisfiability function of the unit-clause formula
\[
 F_{a,b}=
 \left(\bigwedge_{i:a_i=1}z_i\right)\land
 \left(\bigwedge_{i:b_i=1}\neg z_i\right).
\]
The experiment uses the incidence encoding \((a,b)\), not the
general formula encoding of \(T_n\). Conversion to conventional
syntax does not by itself transfer an order-dependent bound to
arbitrary encodings or circuits.

\begin{lemma}[Residual-function width bound \claimid{L02}]
\label{lem:residual}
At a fixed input cut, let \(R\) be the number of distinct functions
of the unread suffix induced by assignments to the read prefix.
Every correct deterministic layered diagram has at least \(R\)
reachable states there. A diagram whose states are residual
functions attains the bound.
\end{lemma}
\begin{proof}
If two prefixes inducing different residual functions reached the
same state, a common suffix would require different answers. A
deterministic diagram cannot produce both. For attainability, take
one state for every residual function at every layer. Each edge
restricts the next variable, defining consistent transitions and
the correct terminal value.
\end{proof}

\begin{proposition}[Grouped-order lower bound \claimid{P03}]
\label{prop:grouped}
For order \(a_1,\ldots,a_m,b_1,\ldots,b_m\), the minimum width
after reading \(a\) is \(2^m\).
\end{proposition}
\begin{proof}
For distinct \(a,a'\), choose \(i\) with \(a_i\ne a'_i\).
Let \(e_i\) have only its \(i\)-th bit set. Then
\(g_m(a,e_i)=1-a_i\) and \(g_m(a',e_i)=1-a'_i\).
All \(2^m\) prefixes induce different residual functions.
Apply Lemma~\ref{lem:residual}.
\end{proof}

\begin{proposition}[Interleaved-order upper bound \claimid{P04}]
\label{prop:interleaved}
For order \(a_1,b_1,\ldots,a_m,b_m\), every layer needs at most
three states. Thus a layered diagram has \(O(m)\) nodes.
\end{proposition}
\begin{proof}
After a completed pair, a prefix either already contains a conflict
or leaves the same disjointness function on the remaining pairs.
There are at most two residual functions. After the first bit of
the next pair, the possibilities are a previous conflict, an
unrestricted next \(b_i\), or the requirement \(b_i=0\).
There are at most three residual functions.
Apply Lemma~\ref{lem:residual} to the \(2m+1\) layers.
\end{proof}

The function also has a linear-size general circuit. The exponential
grouped-order bound concerns a restricted representation. An
explicit diagram needs many nodes in that order; this says nothing
comparable about the length of a program generating the diagram.

\begin{table}[htbp]
\centering
% Generated by experiments/obdd_residuals.py; do not edit.
\begin{tabular}{rrrr}
\toprule
$m$ & Grouped cut & Grouped maximum & Interleaved maximum \\
\midrule
1 & 2 & 2 & 2 \\
2 & 4 & 4 & 3 \\
3 & 8 & 8 & 3 \\
4 & 16 & 16 & 3 \\
5 & 32 & 32 & 3 \\
6 & 64 & 64 & 3 \\
7 & 128 & 128 & 3 \\
\bottomrule
\end{tabular}

\caption{Exhaustive residual-function counts. Grouped cut is the width
after the first \(m\) bits; interleaved maximum includes all layers,
including terminals. Finite observations illustrate the proofs.}
\label{tab:widths}
\end{table}
